\chapter{Introduction}
\label{chap:introduction}

Formal theorem proving has undergone a quiet transformation.
Systems that once required years of expert interaction to close a single non-trivial proof
now close thousands per day --- automatically, verifiably, at scale.
The leading systems are remarkable: Seed-Prover achieves 99.6\% on miniF2F;
BFS-Prover~V2 reaches 95.08\%; Kimina-Prover 92.2\%.
Each of them is a substantial engineering project --- reinforcement learning pipelines,
fine-tuned models, tree search over thousands of candidate tactics,
and verification infrastructure that spans multiple machines.

This study asks a different question.

\section{The gap in the literature}

Every paper that proposes a new search strategy, a new fine-tuning objective,
or a new prompt structure implicitly assumes a baseline.
That baseline is usually described briefly --- ``without our method, pass rate falls to X\%'' ---
and the X is rarely characterised in detail.
We do not know \emph{why} the unimproved agent fails.
We do not know whether it fails because it cannot identify the right tactic,
because it ignores the feedback Lean returns, because it hallucinates lemma names
that do not exist in Mathlib, or simply because it runs out of attempts on a proof
that would have yielded with one more step.

Without a detailed failure taxonomy, every proposed improvement is an educated guess
about which failure mode it addresses.

\section{What this study does}

We characterise the baseline.

The agent is as simple as possible: a single frontier language model ---
NVIDIA Nemotron-3-Ultra-550B, 550~billion parameters, accessed via the NVIDIA NIM API ---
operating in a greedy loop.
At each step it receives the current Lean proof state and the history of what it has tried;
it proposes one tactic; Lean either accepts it, rejects it with an error, or ignores it
because the goal did not change.
No tree search.
No fine-tuning.
No auxiliary tools.
No backtracking.
One tactic at a time until the proof closes or the budget runs out.

We run this agent on 837 real Mathlib theorems, drawn from eleven mathematical domains
and stratified into three difficulty tiers.
Every agent--Lean exchange is recorded in a structured JSONL event log ---
the exact prompt sent, the raw completion, the tactic extracted, the Lean response,
and the step outcome.
That log is the raw material for five questions: what pass rate the agent achieves
per tier, why it fails when it fails, how many steps a successful proof takes,
whether it uses the tactic history it is given, and which tactics it reaches for
versus which ones actually close goals.

\section{Why \emph{this} agent on \emph{these} theorems}

Benchmark suites like miniF2F and PutnamBench are competition mathematics.
They are well-studied, their statements are clean, and their proofs tend to be short.
Real Mathlib theorems are different: they assume a working knowledge of the library,
their proof states can be large, and the right tactic is often one that searches
Mathlib's namespace rather than one that a language model is likely to guess.
A system that scores 99.6\% on miniF2F may behave very differently on the theorems
that practising mathematicians actually leave unproved.

The choice to use a greedy, zero-shot agent is deliberate.
It is the worst reasonable configuration --- no search, no adaptation, no prior training
on the target distribution.
Its failure modes are therefore the \emph{purest} signal: they are not masked by search
finding a workaround, or fine-tuning compressing a latent fix into the weights.
If we understand exactly how this agent fails, we know exactly what each proposed
improvement must overcome.

\section{Contributions}

This study contributes:

\begin{enumerate}
  \item A reproducible evaluation harness for Lean~4 theorem proving,
        model-agnostic and open-source --- the model is selected at runtime with a
        \texttt{--model} flag, and any OpenAI-compatible endpoint works without a
        code change.
  \item 837 Mathlib theorems across three difficulty tiers, with per-step JSONL logs
        for every agent--Lean exchange.
  \item Pass rates for Nemotron-3-Ultra-550B and Groq~\texttt{gpt-oss-20b} on the Easy
        tier, reported both as a headline figure over all theorems attempted and as a
        reachable-only figure over those that actually reached a working model ---
        a distinction the runs made unavoidable, and one that changes the number
        substantially. Medium and Hard remain to be run.
  \item The first published failure-mode taxonomy for LLM-based Lean agents:
        five mutually exclusive modes derived from the event logs, with examples.
  \item Evidence on history utilisation: whether the agent reads and acts on
        the tactic history it receives.
\end{enumerate}

The result is the characterisation that every subsequent paper in this space can use
as its reference point.
